\label{chap:cobertura-taxonomia-numeros}
\index{HMT number!taxonomy}
\index{Hensel transducer!complete displacement}
\index{rational number!periodicity}

\cref{chap:numero-medicion} defined an HMT number as a compatible section
of the inverse survival system, and not as a sequence of
digits devoid of genealogy. There remained, however, two questions that
should be separated precisely. The first is a question of
\emph{representation}: which real values can the
3-adic transducer represent before the particular closure filters are imposed? The second is
a question of \emph{selection}: which sections of the raw space are
selected by APP, TRIT, and TPK when their compatibilities are required
simultaneously? This section resolves the first, formulates the second as
an explicit domain restriction, and extracts from the base-one-thousand chart a
complete taxonomy of visible values. The foundational tools are
\(p\)-adic arithmetic, symbolic dynamics, and the positional characterization
of the rational numbers \cite{Koblitz1984,LindMarcus1995,HardyWright2008}.

The distinction is structural. A language capable of representing every tape does not, by that fact alone, select the tape of \(\pi\), \(e\), \(\varphi\), or \(\alpha\). Conversely, a selection law cannot be evaluated unless the space on which it acts has first been constructed. Representation and selection are two successive theorems, not two names for the same step.

\section{Unfiltered Hensel transducer and full shift on the ternary fibers}

Let \(p\) be a prime, let \(d\geq1\), and let
\(A\in\operatorname{GL}_d(\mathbb Z)\). For a row vector
\(x_t\in\mathbb Z_p^d\), we define
\begin{equation}
 y_t=x_tA,
 \qquad
 u_t=y_t\bmod p\in\mathbb F_p^d,
 \qquad
 x_{t+1}=\frac{y_t-u_t}{p}.
 \label{eq:transductor-hensel-general}
\end{equation}
The block \(u_t\) is the visible emission, and \(x_{t+1}\) is the memory that
remains after that digit is removed.

\begin{theorem}[Hensel Conjugation with Full Shift]
\label{thm:shift-hensel-completo}
For every \(T\geq1\), the map
\[
 \Psi_T:(\mathbb Z/p^T\mathbb Z)^d
 \longrightarrow(\mathbb F_p^d)^T,
 \qquad
 x_0\longmapsto(u_0,\ldots,u_{T-1}),
\]
is bijective. Its inverse is
\begin{equation}
 x_0\equiv
 \sum_{t=0}^{T-1}p^t u_tA^{-(t+1)}
 \pmod {p^T}.
 \label{eq:inversa-hensel-finita}
\end{equation}
Taking the inverse limit yields a homeomorphism.
\[
 \Psi:\mathbb Z_p^d\xrightarrow{\ \sim\ }
 (\mathbb F_p^d)^{\mathbb N}.
\]
If \(S_A\) denotes the update \(x_t\mapsto x_{t+1}\) and
\(\sigma\) removes the first block from a tape, then
\[
 \Psi\circ S_A=\sigma\circ\Psi.
\]
\end{theorem}

\begin{proof}
From \cref{eq:transductor-hensel-general} is obtained
\[
 x_t=(u_t+px_{t+1})A^{-1}.
\]
Iteration of this identity gives
\[
 x_0=
 \sum_{t=0}^{T-1}p^t u_tA^{-(t+1)}
 +p^Tx_TA^{-T}.
\]
The final summand vanishes modulo \(p^T\), which proves that a word of
\(T\) blocks determines at most one initial class. Conversely, given
any word, \(x_T=0\) is fixed and reconstruction proceeds backward by means of
\(x_t=(u_t+px_{t+1})A^{-1}\). The equality
\(x_tA=u_t+px_{t+1}\) verifies that the resulting class emits the prescribed
word. Therefore, \(\Psi_T\) is bijective.

The bijections are compatible with reduction modulo \(p^T\). Their inverse
limits produce the homeomorphism \(\Psi\); equivalently, the series of
\cref{eq:inversa-hensel-finita} converges in \(\mathbb Z_p^d\). Applying
\(S_A\) removes exactly the block \(u_0\), and therefore corresponds to
applying \(\sigma\) on the strip.
\end{proof}

In the HMT instance, \(p=3\), \(d=6\), and the published matrix has
determinant one. Hence
\begin{equation}
 \mathbb Z_3^6\simeq(\mathbb F_3^6)^{\mathbb N},
 \qquad |\mathbb F_3^6|=3^6=729.
 \label{eq:hmt-shift-729}
\end{equation}
Each tape of six-trit blocks has a unique initial 3-adic memory.
The matrix determines the coordinatization and the reading dynamics; the
universality of the alphabet proves that it excludes no tape.

\begin{corollary}[Entropy and Haar measure]
The raw Hensel update is conjugate to the full one-sided shift on
\(729\) symbols. Its topological entropy is \(6\log3\). The Haar measure on
\(\mathbb Z_3^6\) is transported to the uniform product measure on the
blocks.
\end{corollary}

\begin{proof}
There are \(729^T=3^{6T}\) cylinders of length \(T\), and each class modulo
\(3^T\) has Haar measure \(3^{-6T}\). Conjugation by
\cref{thm:shift-hensel-completo} transports both properties to the
shift.
\end{proof}

The probabilistic conclusion concerns the raw space endowed with Haar measure. It
does not assert that the published sections of the constants were sampled
according to that measure or that they are algorithmically random.

\section{Balanced ternary chart of the integers}

The raw 3-adic chart and the Archimedean chart do not exhaust the numerical types
of the HMT classification. For the integral operator, the empty word is used together with the
reduced words
\[
 \mathscr W_{\mathrm{bal}}
 :=\{\varnothing\}\cup
 \{a_1\cdots a_m:
 a_j\in\{-1,0,+1\},\ a_1\ne0\}.
\]
We define recursively
\[
 b(\varnothing)=0,
 \qquad
 b(wa)=3b(w)+a.
\]

\begin{theorem}[Unique balanced ternary representation]
\label{pf84:C03}
The evaluation
\[
 b:\mathscr W_{\mathrm{bal}}\longrightarrow\ZZ
\]
is bijective. Words whose first trit is \(+1\) represent the
positive integers, and those beginning with \(-1\) represent the negative integers.
\end{theorem}

\begin{proof}
For \(n\ne0\), the ultimo trit \(a\) is the representante unique of
\(n\bmod3\) in \(\{-1,0,+1\}\). The padre
\[
 \frac{n-a}{3}
\]
has strictly smaller absolute value. The iteration terminates at \(0\), so
that every integer has a finite expansion. The uniqueness of the balanced
residue at each step proves the uniqueness of the word. The sign is fixed
by the first nonzero trit.
\end{proof}

\begin{remark}
This bijection is an exact arithmetic chart. It does not identify all
balanced words with surviving TPK sections, nor does it replace the
genealogy, carry, or orientation of an enriched route.
\end{remark}

\section{Surjectivity of the unfiltered Archimedean chart onto \(\mathbb R\)}

Let us concatenate the coordinates of the blocks
\(u_t\in\mathbb F_3^6\) to obtain a tape
\((a_n)_{n\geq1}\in\{0,1,2\}^{\mathbb N}\). Its ternary evaluation is
\[
 \operatorname{ev}_3((a_n))=
 \sum_{n\geq1}\frac{a_n}{3^n}\in[0,1].
\]
We adopt the canonical convention that excludes a representation with an
eventually constant tail of twos only when the same point has another expansion
in the chart. The endpoint \(1\) therefore retains its unique fractional
representation \(0.222\ldots{}_3\).

\begin{theorem}[Archimedean raw coverage]
\label{thm:cobertura-arquimediana-bruta}
The map
\[
 P_{\rm raw}:\mathbb Z_3^6
 \xrightarrow{\Psi}(\mathbb F_3^6)^{\mathbb N}
 \xrightarrow{\operatorname{flat}}\{0,1,2\}^{\mathbb N}
 \xrightarrow{\operatorname{ev}_3}[0,1]
\]
is continuous and surjective. Therefore,
\[
 \mathbb Z_3^6/\!\sim_{P_{\rm raw}}
 \ \cong\ [0,1]
\]
as the value space, where \(x\sim_{P_{\rm raw}}y\) if their evaluations coincide.
\end{theorem}

\begin{proof}
Every ternary digit belongs to a unique consecutive block of six digits, so \(\operatorname{flat}\) is a bijection of tape spaces. Every \(r\in[0,1]\) has a ternary expansion; the \cref{thm:shift-hensel-completo} provides the unique 3-adic memory that emits its blocks. Continuity is verified by cylinders: fixing \(T\) blocks fixes \(6T\) digits and confines the value to an interval of diameter \(3^{-6T}\). Identification of the quotient with the image is the canonical factorization of a surjective map by its fibers.
\end{proof}

A single 3-adic chart is compact and therefore cannot continuously cover
an unbounded line. The correct globalization uses a countable family
of charts.

\begin{corollary}[Cover Archimedean countable of \(\mathbb R\)]
\label{cor:atlas-hmt-reales}
Let
\[
 \widehat X_{\rm raw}=\mathbb Z\times\mathbb Z_3^6,
 \qquad
 \widehat P_{\rm raw}(m,x)=m+P_{\rm raw}(x).
\]
Then \(\widehat P_{\rm raw}\) is surjective onto \(\mathbb R\).
Consequently, every real number admits at least one raw Hensel genealogy and
\[
 \widehat X_{\rm raw}/\!\sim_{\widehat P_{\rm raw}}
 \ \cong\ \mathbb R
\]
as a set of values.
\end{corollary}

\begin{proof}
For \(r\in\mathbb R\), tome \(m=\lfloor r\rfloor\) y aplique the
\cref{thm:cobertura-arquimediana-bruta} to \(r-m\in[0,1)\).
\end{proof}

This result makes precise the assertion that \(\mathbb R\) is the Archimedean
quotient image of a genealogical structure: the projection forgets 3-adic memory and
preserves the evaluated point. The assertion proved here concerns the raw completion
\emph{raw}. If \(X_{\rm surv}\subseteq\mathbb Z_3^6\) denotes the subspace
that also satisfies a concrete family of APP--TRIT--TPK closures, then
\[
 P_{\rm raw}(X_{\rm surv})=[0,1]
\]
is an additional theorem. The transducer guarantees full capacity; survival
determines what portion of that capacity the restricted system
realizes.

Nor is it deduced here that the native sum and product of TPK descend to
the entire global quotient. The usual arithmetic may be transported from
\(\mathbb R\), but proving that it coincides with operations constructed before
evaluation requires the corresponding descent theorem.

In this construction, the cylinders are not nonzero infinitesimals of
\(\mathbb R\). Every finite-depth cylinder has positive diameter,
and the sequence of diameters converges to zero. Two sections that agree up to
depth \(T\) and differ thereafter form a \emph{tail
perturbation}; their Archimedean difference is bounded by the diameter of the common
cylinder. This is the exact formulation of the ``cylinders tending to zero''
within the ordinary real line.

\section{Positional chart in base \(10^3\) and classification of the published sections}

The cubic completion
\(1000=729+243+27+1\) makes it natural to group the decimal digits into blocks
\(b_n\in\{0,\ldots,999\}\). For a tape of blocks, we define
\[
 \operatorname{ev}_{1000}(b_1,b_2,\ldots)
 =\sum_{n\geq1}\frac{b_n}{1000^n}.
\]
The canonical expansion excludes an eventually constant tail of blocks
\(999\). This convention resolves the duplication of cylinder endpoints.

\begin{theorem}[Visible taxonomy based on thousand]
\label{thm:taxonomia-base-mil}
Let \(x\in[0,1)\) be given, and let \((b_n)\) be its canonical base-one-thousand expansion.
Then:
\begin{enumerate}
\item \(x\) has a terminating expansion if and only if, when written as
an irreducible fraction \(a/q\), one has \(q\mid1000^N\) for some \(N\);
equivalently, the only primes dividing \(q\) are \(2\) and \(5\).
\item \(x\) is rational if and only if \((b_n)\) is eventually periodic.
\item \(x\) is irrational if and only if \((b_n)\) is not eventually
periodic.
\end{enumerate}
\end{theorem}

\begin{proof}
An expansion terminating at depth \(N\) has the form
\(D/1000^N\); upon reduction, its denominator contains only the primes \(2\) and
\(5\). Conversely, if \(q\mid1000^N\), then \(a/q\) is written with
\(N\) blocks.

If, after a prefix of length \(r\), a period of length
\(s\) is repeated whose base-thousand integer is \(P\), the tail is a geometric series:
\[
 \frac{P}{1000^r(1000^s-1)},
\]
and the value is rational. Conversely, when \(a/q\) is expanded, the remainders of
successive division can assume only \(q\) values. A zero remainder produces a
terminating expansion; repetition of a remainder produces a period. The third
statement is the exact negation of the second.
\end{proof}

The word \emph{fraction} denotes a presentation \(a/q\) of a rational
number, not an additional numerical species. Nor should periodicity and
algebraicity be conflated. A real number is algebraic when it is a zero of
some nonzero polynomial over \(\mathbb Q[X]\); it is transcendental otherwise. This
second classification is independent of the base used to write its digits.
For example,
\[
 \varphi^2-\varphi-1=0
\]
makes \(\varphi\) an algebraic irrational, whereas the transcendence
of \(e\) and \(\pi\) proceeds from classical theorems external to HMT\@.

\section{Concatenation, nonadic reduction, and carry}

Compatibility between the cubic chart and the digital root is an arithmetic
identity, not a visual resemblance. If \(u_j\in\{0,\ldots,999\}\) and
\[
 D_n=1000D_{n-1}+u_n,
 \qquad D_0=0,
\]
then \(D_n\) is the integer obtained by concatenating the first \(n\)
blocks.

\begin{theorem}[Nonadic congruence of control of the cubic chart]
\label{thm:checksum-nonadico-base-mil}
For every \(n\geq1\),
\[
 D_n\equiv\sum_{j=1}^n u_j\pmod9.
\]
Therefore, concatenation in base ten, reduction modulo nine, and digital root are compatible.
\end{theorem}

\begin{proof}
Since \(1000\equiv1\pmod9\), the recurrence gives
\(D_n\equiv D_{n-1}+u_n\pmod9\). Iteration from \(D_0=0\) produces the
formula.
\end{proof}

In the dodecaphase closure of \(\alpha\), the blockwise carry equation
has the form
\[
 P_m+E_m-\Phi_m-K_m+c_{m+1}=A_m+1000c_m.
\]
Reduced modulo nine,
\[
 A_m\equiv
 P_m+E_m-\Phi_m-K_m+c_{m+1}-c_m
 \pmod9.
\]
If the extreme carries vanish, the telescoping sum eliminates the intermediate
memory. The global residue of the output then coincides with the residue
of the combination of inputs. This law explains why the nonadic record
controls decimal concatenation; by itself, it does not select the blocks that
must be concatenated.

\section{Nonadic elevated coordinate and phase return}

Reduction modulo nine must not be confused with the identification of the
integers that possess the same residue. Let
\[
 I_9=\{1,2,\ldots,9\}.
\]
For each positive integer \(n\), we define its \emph{nonadic phase} and its
\emph{revolution index} by
\begin{equation}
 \rho_9(n)=1+((n-1)\bmod9),
 \qquad
 \kappa_9(n)=\frac{n-\rho_9(n)}9.
 \label{eq:coordenada-nonadica-levantada}
\end{equation}
The first coordinate is the digital root written in the alphabet
\(I_9\), so that the zero residue class is represented by phase \(9\).
The second records how many complete turns have preceded the
visible phase.

\begin{theorem}[Nonadic decomposition with memory]
\label{thm:descomposicion-nonadica-memoria}
The map
\[
 \Theta_9:\mathbb N_{>0}\longrightarrow\mathbb N_0\times I_9,
 \qquad
 n\longmapsto\bigl(\kappa_9(n),\rho_9(n)\bigr),
\]
is bijective, with inverse
\[
 \Theta_9^{-1}(q,r)=9q+r.
\]
In these coordinates, the arithmetic successor is conjugated with the dynamics.
\begin{equation}
 \mathcal S_9(q,r)=
 \begin{cases}
  (q,r+1),&1\leq r<9,\\
  (q+1,1),&r=9.
 \end{cases}
 \label{eq:sucesor-nonadico-levantado}
\end{equation}
In particular, the visible transition \(9\to1\) is a phase return accompanied
by the irreversible increment \(q\to q+1\); it is not an equality between the
integers \(9\) and \(10\).
\end{theorem}

\begin{proof}
Euclidean division of \(n-1\) by \(9\) uniquely provides
\(n-1=9q+s\), with \(q\in\mathbb N_0\) and \(0\leq s<9\). Taking
\(r=s+1\) yields \(n=9q+r\), with \(r\in I_9\), which proves simultaneously
existence, uniqueness, and the inverse formula. If \(r<9\), adding one unit only
increments \(r\). If \(r=9\), the equality
\(9q+9+1=9(q+1)+1\) gives the second branch of
\cref{eq:sucesor-nonadico-levantado}.
\end{proof}

This elevation separates three notions that the isolated figure mixes:
\begin{enumerate}
\item the unit of advancement, realized by the successor;
\item the visible phase \(r\), which traverses the nonadic cycle;
\item the accumulated memory \(q\), which distinguishes returns of the same phase.
\end{enumerate}
The projection \((q,r)\mapsto r\) forgets the number of turns and satisfies
\[
 \rho_9(n+9)=\rho_9(n),
 \qquad
 \kappa_9(n+9)=\kappa_9(n)+1.
\]
This is the elementary arithmetic model of the mechanism that, in the enriched
TPK state, additionally preserves Hensel quotients, orientation, signature, and
boundary datum. A return of the visible word can thus coexist with a
distinct genealogy.

Zero does not belong to \(I_9\): it is added as a separate base point before \(\Theta_9\) is applied. For negative integers, the positive modulus of the decomposition is preserved and an orientation \(\varepsilon\in\{-1,+1\}\) is appended. Consequently, neither the sign nor zero is deduced from the residual phase. This typing prevents a cyclic chart from being converted into a false arithmetic identity.

\begin{theorem}[Oriented Integer and Genealogical Sign]
\label{thm:entero-orientado-signo-genealogico}
The map
\[
\Theta_9^\pm:
\mathbb Z\setminus\{0\}
\longrightarrow
\{-1,+1\}\times\mathbb N_0\times I_9,
\]
\[
n\longmapsto
\bigl(\operatorname{sgn}(n),
      \kappa_9(|n|),
      \rho_9(|n|)\bigr)
\]
is bijective, with inverse
\[
(\varepsilon,q,r)\longmapsto\varepsilon(9q+r).
\]
Aritmetical negation acts by
\[
(\varepsilon,q,r)\longmapsto(-\varepsilon,q,r):
\]
reverses the orientation and preserves magnitude, phase, and turn memory.
\end{theorem}

\begin{proof}
The \cref{thm:descomposicion-nonadica-memoria} uniquely provides
\(|n|=9q+r\), with \(q\in\mathbb N_0\) and \(r\in I_9\). The sign of a nonzero
integer is likewise unique. The inverse formula recomposes \(n\), and replacing
\(n\) with \(-n\) changes only \(\varepsilon\).
\end{proof}

Thus, a negative is neither a special residual phase nor an
``impossible'' quantity. It is the same nonadic magnitude transported with the opposite
genealogical orientation. This structure will be prolonged in Dimensional Mechanics
by the particle--antiparticle anti-section.

Iterating \(\kappa_9\) produces a finite tower for every ordinary integer; in the inverse completion, a compatible sequence of quotients may be infinite. The distinction between the two cases anticipates the HMT taxonomy: the visible value is determined by the positional expansion, whereas the genealogical type is determined by the complete evolution of its memories.

\section{2 classification axes: visible value and transport genealogy}

The preceding taxonomy classifies the projected value. HMT adds a second
axis that ordinary decimal notation discards. For a section
\(x\in\Omega_\infty^{\rm Arch}\), we distinguish:
\[
 \begin{array}{rcl}
 \text{visible type}&=&
 \text{terminating, periodic or nonperiodic; algebraic or transcendental},\\
 \text{genealogical type}&=&
 \text{carry class, orientation, closure and monodromy of }x.
 \end{array}
\]
Two sections may share the same real value and differ on the second axis. A
rational real has an eventually periodic visible expansion, but a
redundant enriched presentation may preserve nonperiodic internal memory
that evaluation forgets. Periodicity of the \emph{value} always
refers to its canonical expansion; periodicity of the
\emph{genealogy} refers to the TPK enriched state.

This separation provides a concise definition:
\begin{equation}
 \boxed{
 \text{HMT number}
 =
 \text{compatible section with memory};
 \qquad
 \text{real value}
 =
 \text{its evaluation class}.}
 \label{eq:numero-hmt-y-valor-real}
\end{equation}
Measurement in this chart consists in coupling the section to the reader, fixing a cylinder, and
conditioning the compatible prolongations. Increasing precision means
narrowing the cylinder, not consuming the genealogy that remains in the fibre.

\section{Memorization of arbitrary depth in finite emission systems}

The visible recurrence of nine phases can coexist with an irrational
output only if the complete state preserves information that does not reduce to
those nine labels. The following result makes this necessity exact.

\begin{theorem}[Characterization via finite emitters]
\label{thm:emisor-finito-racional}
Let \(Q\) be a finite set, \(F:Q\to Q\) a deterministic transition,
\(o:Q\to\{0,\ldots,999\}\) an output, and \(q_0\in Q\). The tape
\(b_n=o(F^{n-1}(q_0))\) is eventually periodic and its evaluation is rational.
Conversely, every eventually periodic tape admits a finite deterministic
emitter.
\end{theorem}

\begin{proof}
Among \(|Q|+1\) consecutive states, two are equal. From the first
repetition, determinism forces the entire future to repeat, so the
tape is eventually periodic. The \cref{thm:taxonomia-base-mil} gives
rationality. For the converse, it suffices to construct a chain of states for the
prefix and a cycle for the period.
\end{proof}

Therefore, a law generating exactly \(e\), \(\pi\), or any other
irrational number cannot consist solely of a deterministic finite-state
automaton that repeats without additional memory. It may use an inverse
limit, 3-adic memory, increasing depth, or an
aperiodic input; all these mechanisms provide unbounded effective state.
In HMT, visible-phase return
\(g_{\rm vis}(n+9)=g_{\rm vis}(n)\) does not force return of the carry,
orientation, or boundary. That difference is precisely where
an aperiodic expansion may reside.

\section{The local square of \texorpdfstring{\(e\)}{e} and the memory break}

For a decimal word \(a_1\cdots a_n\), we shall call a factorization a \emph{local square}
of length \(\ell\) at position \(i\)
\[
 a_i\cdots a_{i+2\ell-1}=vv,
 \qquad |v|=\ell.
\]
This notion belongs to word combinatorics. It does not assert that the
infinite tail has a period.

\begin{theorem}[Finite characterization of the square \(1828^2\)]
\label{thm:cuadrado-local-e}
Consider the \(243\) catalog words and the fixed reader
\(w_6\mapsto w_{30}\mapsto d_1d_2d_3d_4\), with four triads in base one thousand.
There is a unique word whose reading contains a square of length four
beginning at the second decimal digit. It is
\[
 w_e=201101,
 \qquad
 718281828459=7\,\underbrace{1828\,1828}_{(1828)^2}\,459.
\]
The repetition is not prolonged to a third copy: the digit that would require that
continuation would be \(1\), whereas the effective digit is \(4\). In the whole
catalogue there exists only one other square of length four,
\[
 575157518472=\underbrace{5751\,5751}_{\text{local square}}\,8472,
\]
which begins at the first digit.
\end{theorem}

\begin{proof}
The exhaustive enumeration evaluates the \(243\) words using integer
arithmetic. The square profile for lengths \(1,\ldots,6\) contains
\[
 240,\ 21,\ 3,\ 2,\ 0,\ 0
\]
occurrences, distributed among
\[
 153,\ 19,\ 3,\ 2,\ 0,\ 0
\]
words. The two occurrences of length four are exactly those in the
statement. Two independent integer enumerations reproduce the same census,
and their complete outputs coincide term by term.
\end{proof}

The word \(w_e\) also has the lifting chain
\[
201101\mid121221\mid102011\mid012222\mid102011.
\]
The third and fifth visible blocks coincide. The states that transport them,
however, do not coincide. Before both emissions, their reductions
modulo \(27\) are
\[
(25,11,7,17,8,16),
\qquad
(7,8,4,23,26,7),
\]
and the subsequent Hensel quotients are
\[
(6,13,9,2,13,1),
\qquad
(15,0,3,19,23,25).
\]
Therefore, the return \(102011\) is a return of the visible projection, not
of the enriched state. This is the exact formulation of the local repetition
that is broken: the quotient preserves a genealogy that the residual block does not
record.

Consequently, the beginning
\[
e=2.718281828459\ldots
\]
does not exhibit a ``semiperiodicity,'' but rather a local word-square with
visible-phase return and memory advance without state periodicity. The result is exhaustive within the
fixed reader; its content is combinatorial and dynamical, not an inference of
irrationality from twelve digits.

\section[2 finite measures of the emitter 468 a 243]{2 finite measures of the \texorpdfstring{\(468\to243\)}{468→243} emitter: uniform frequency and structural priority}

The catalogue publico distinguishes \(468\) emissions y \(104\,976\) semillas.
Be
\[
 q:E_{468}\longrightarrow W_{243}
\]
the map that forgets the emission signature and preserves the six-trit
word. There are two natural measures because there are two finite counting universes:
\[
 \mu_{468}(w)=\frac{|q^{-1}(w)|}{468},
 \qquad
 \mu_{\rm seed}(w)=
 \frac{\sum_{e\in q^{-1}(w)}m(e)}{104\,976},
\]
where \(m(e)\) is the multiplicity of semillas of the emision \(e\).

\begin{theorem}[Canonical measures relative to the counted universe]
\label{thm:medidas-emisor-468}
The uniform measure on \(E_{468}\) is the unique probability invariant under
all permutations of its emissions; its image measure is \(\mu_{468}\).
The uniform measure on the \(104\,976\) seeds induces
\(\mu_{\rm seed}\). Their exact histograms are
\[
\begin{array}{c|rrrrrr}
|q^{-1}(w)|&1&2&3&4&5&6\\ \hline
\#\{w\}&132&66&18&3&6&18
\end{array}
\]
and
\[
\begin{array}{c|rrrrrrrrr}
\text{seed weight}
&144&288&432&576&864&1008&1440&1944&2088\\ \hline
\#\{w\}
&120&54&18&9&15&6&3&12&6.
\end{array}
\]
For the three words distinguished by the reader, one obtains
\[
\begin{array}{c|cc}
&\mu_{468}&\mu_{\rm seed}\\ \hline
\pi&5/468&7/729\\
e&1/468&1/729\\
\varphi&1/468&1/729.
\end{array}
\]
\end{theorem}

\begin{proof}
Invariance under the symmetric group assigns equal masses to all
points of \(E_{468}\); normalization fixes them at \(1/468\). The formula for
the pushforward measure follows by summing over each fiber. The same argument applied to the
seed set yields the second measure. The histograms are obtained
by exhaustive enumeration of the finite catalogue of \(468\) emissions and \(243\)
words; their sums verify
\[
132+66+18+3+6+18=243,
\]
\[
132+2\cdot66+3\cdot18+4\cdot3+5\cdot6+6\cdot18=468
\]
and the sum total of the weights is \(104\,976\).
\end{proof}

The nonuniformity of the pushforward measure is a finite theorem: distinct words have different numbers of realizations. Conversely, no normalizable uniform probability exists on all of \(\mathbb R\), and these measures do not intrinsically order all the reals. The word for \(\pi\) belongs to one of the six fibers of size five and weight \(1008\); it is not the unique maximum of either histogram. Its proved singularity arises from the composite structure of its five regions, from its oriented partition \(3+2\), and from its subsequent compatibilities, not from declaring an absolute probability on the continuum a priori.

\section{Selector orbital finite preceding the Archimedean reader}
\label{sec:selector-orbital-constantes}

The preceding measures reveal multiplicities, but an isolated six-trit word
still depends on an orientation. The catalog has an internal symmetry that
allows unoriented classes to be selected first and a representative to be fixed afterward.
Write a word as \(abc\mid def\) and define
\[
 r(abc\mid def)=cab\mid efd,
 \qquad
 s(abc\mid def)=def\mid abc.
\]
These permutations also act on the six coordinates of the datum
\(U_6\). In the complete catalog, they satisfy
\[
 r^3=1,
 \qquad s^2=1,
 \qquad srs=r^{-1},
\]
and hence they generate an action denoted by \(D_3^{\rm cat}\). The action uses
the fixed TRIT labeling
\[
 \Psi(U_6)=(-U_6\bmod3)
\]
and the published partition of the six coordinates into two triads.

For a word \(w\), let \(f(w)=|q^{-1}(w)|\) be the number of emissions,
\(m(w)\) its seed mass, and
\[
 \nu(w)=(n_0(w),n_1(w),n_2(w))
\]
its ordered census of trits. Also let \(d(e)\) be the additive diagonal class of
an emission and \(o(e)\in\{\mathrm{ES},\mathrm{NO}\}\) its published
orientation.

\begin{theorem}[Internal orbital selection of the three words]
\label{thm:selector-orbital-constantes}
The action \(D_3^{\rm cat}\) preserves the set of \(468\) emissions, the
projection \(q:E_{468}\to W_{243}\), and the seed multiplicity. Its action on
\(W_{243}\) has exactly \(43\) orbits: \(38\) of size six and
five of size three. The following uniqueness statements hold within that
quotient.

\begin{enumerate}
\item There exists a unique orbit \(\mathcal O_\pi\) of size six such that, for
every \(w\in\mathcal O_\pi\),
\[
 f(w)=5,
 \qquad m(w)=1008,
 \qquad
 \{m(e):e\in q^{-1}(w)\}=\{432,144,144,144,144\}.
\]
It is
\[
\mathcal O_\pi=
\{001112,010211,100121,112001,121100,211010\},
\]
and all its words have census \(\nu=(2,3,1)\).

\item Consider the orbits of size six whose words have a single
emission, of multiplicity \(144\), and whose joint diagonal support is
\(\{0,3,6\}\). There are exactly six. Among them there is a unique orbit
with the same census \((2,3,1)\) of \(\mathcal O_\pi\):
\[
\mathcal O_e=
\{011120,012110,101201,110012,120011,201101\}.
\]

\item In the same sector there exists a unique balanced orbit, that is, with
\(\nu=(2,2,2)\):
\[
\mathcal O_\varphi=
\{002112,020211,112002,121200,200121,211020\}.
\]
\end{enumerate}

Finally, the guiding clause
\begin{equation}
 \exists e\in q^{-1}(w):
 \qquad d(e)=6,
 \qquad o(e)=\mathrm{ES}
 \label{eq:gauge-orbital-constantes}
\end{equation}
select a single representative from each orbit and produces, respectively,
\[
 \boxed{010211,\qquad201101,\qquad121200.}
\]
\end{theorem}

\begin{proof}
The group relations are verified directly as identities of permutations of six
coordinates. Enumeration of the \(468\) rows
verifies that \(r\) and \(s\) send every \(U_6\) to another row of the catalogue with
the same multiplicity, and that ternary projection commutes with both actions.
Therefore, the orbits can be computed on the \(243\) words without losing the
fibers.

The exhaustive partition gives the size histogram.
\[
 38\cdot6+5\cdot3=243.
\]
Filtering those \(43\) orbits by the first predicate leaves exactly one;
filtering by singleton fiber, minimum multiplicity, and diagonal support
\(\{0,3,6\}\) leaves six. The census \((2,3,1)\) leaves a single orbit in the
first case, and the balanced census \((2,2,2)\) leaves a single orbit in the second.
Finally, the search for a lift with \(d=6\) and orientation
\(\mathrm{ES}\) leaves one word in each of the three orbits. All
filtrations are performed on structural descriptors of the catalogue before consulting the
decimal reader.
\end{proof}

The predicates of the theorem may be combined into a selection invariant.
For an orbit \(\mathcal O\), we define its \emph{orbital profile}
\[
 \Sigma_{\rm cat}(\mathcal O)=
 \bigl(
 |\mathcal O|,\ f,\ m,\ \mathcal M_{\rm fib},\
 \mathcal D,\ \nu
 \bigr),
\]
where \(f\) is the fibre size, \(m\) the seed mass,
\(\mathcal M_{\rm fib}\) the multiset of multiplicities of their
lifts, \(\mathcal D\) the diagonal support when used, and \(\nu\)
the trit census. The three distinguished classes occupy different positions
in this invariant space:
\[
\begin{array}{c|c|c|c|c|c}
\mathcal O&|\mathcal O|&f&m&\mathcal M_{\rm fib}&(\mathcal D,\nu)\\ \hline
\mathcal O_\pi&6&5&1008&\{432,144,144,144,144\}&(-,(2,3,1))\\
\mathcal O_e&6&1&144&\{144\}&(\{0,3,6\},(2,3,1))\\
\mathcal O_\varphi&6&1&144&\{144\}&(\{0,3,6\},(2,2,2)).
\end{array}
\]
The demonstrated singularity does not mean that these numbers receive an
absolute probability over \(\mathbb R\). It means that, within the finite universe
produced by the emitter and before evaluating digits, their classes are
separable by internal invariants and that the class of \(\pi\) has a
regional fiber of a different type. This is the mathematical formulation of a
structural priority in HMT\@.

The negative control is informative. If the diagonal support \(\{0,3,6\}\) is omitted, there exist four balanced singleton orbits and \(\mathcal O_\varphi\) ceases to be unique. Before the gauge is fixed, each orbit admits six possible orientations. Once it is fixed, the word of \(\pi\) preserves three positive lifts of \(U_6\), rather than a single region. Thus, the theorem selects the oriented word while respecting the geometric multiplicity of the five regions.

The non-uniformity is reinforced when passing to orbits:
\[
\begin{array}{c|cc}
&\text{emissions}&\text{seeds}\\ \hline
\mathcal O_\pi&5/78&14/243\\
\mathcal O_e&1/78&2/243\\
\mathcal O_\varphi&1/78&2/243.
\end{array}
\]
These quotients remain masses in the finite catalog, not probabilities
assigned to the evaluated real numbers.

The causal advance is precise. The three words cease to be named inputs: they are recovered from the symmetry and finite invariants of APP--TPK, relative to the action \(D_3^{\rm cat}\), the TRIT labeling, the distribution \(3+3\), and the orienting gauge \cref{eq:gauge-orbital-constantes}. The proof uses neither \(L_0\), \(L_1\), decimal digits, \(\alpha\), CODATA, nor the dodecaphasic state. The absolute uniqueness of that action among all possible automorphisms, the deduction of the gauge rather than its declaration, the canonicity of the reader \(L_0/L_1\) independently of the target value, and the coinductive emission of the infinite expansions remain subsequent claims.

\section{Position of \texorpdfstring{\(\pi,e,\varphi\) and \(\alpha\)}{pi, e, phi, and alpha} in the classification}

The \cref{thm:shift-hensel-completo} explains why any decimal prefix
can be re-emitted from a 3-adic state: the raw capacity is
universal. Therefore, the mere back-and-forth operation on a tape selects no
constant. The distinctive content of the published construction lies
in the objects preceding and following that re-emission:
\begin{enumerate}
\item the finite fibers and the five distinct regions of \(\pi\);
\item the typed Archimedean reader that associates words, cylinders, and blocks;
\item the dodecafase signed and the closure reversible
\[
U_{12}^{\rm sgn}
\longleftrightarrow K
\longleftrightarrow(D_3K,D_4K,Q)
\longleftrightarrow\alpha_{\rm HMT};
\]
\item the contrast criteria that distinguish an equivalence of charts
from a selector preceding evaluation.
\end{enumerate}

The twelve published decimal triples of \(\pi,e,\varphi\) are exact finite
coordinates of the enriched reader. The printed string of thirty-six
digits is the rational dodecaphasic projection
\(\alpha_{12}=\pi_{12}(\alpha_{\rm HMT})\):
\[
\frac{7297352569283800997285105472380663}{10^{36}}.
\]
The exact identity in this window is
\[
\alpha_{12}
:=\pi_{12}(\alpha_{\rm HMT})
=\frac{7297352569283800997285105472380663}{10^{36}}.
\]
The compatible recurrence at arbitrary depth constructs the infinite
section
\[
 \alpha_{\rm HMT}:=\varprojlim_N A^{(N)},
\]
and each rational window is an exact projection of that same object. The
kernel \(E_{21/22}\) constructs, in another chart, the finite simple positive
root \(\alpha_E\); the graded transport of the TPK extends that kernel into
a compatible family whose limit root is \(\alpha_B\). If
\(\alpha_A:=\varprojlim_N\rho_N(\alpha_{12})\), the truncation squares
of the two families commute and determine the same surviving
cylinder at each depth. The exact two-way theorem therefore establishes
\[
 \boxed{\alpha_A=\alpha_B=\alpha_{\rm HMT}}.
\]
The equality of \(\operatorname{Tr}_9(\alpha_E^{-1})\) and
\(\operatorname{Tr}_9(\alpha_{12}^{-1})\) is the first finite certificate of
that identity; it neither defines it, limits its scope nor selects the
extensions. The typed projection keeps the coinductive
section, its rational window and the analytic root distinct without separating their identity
in the limit.

The square root of \(-1\) belongs to another chart. In the
\cref{chap:algebras-cuadraticas-orden-nueve}, one constructs
\[
\mathbb F_9=\mathbb F_3[\iota]/(\iota^2+1),
\qquad \iota^2=-1,
\]
and the Witt matrix is proved to realize that quadratic structure. This is
an exact finite algebraic extension, not a visible category of real
expansions.

\section{Unfiltered Hensel chart, survival, and global boundary}

The following controls are falsifiers of the preceding finite results:
\begin{enumerate}
\item that two states distinct modulo \(3^T\) emit the same word of
\(T\) blocks, or that there is a word without a preimage;
\item failure of the explicit inverse of
\cref{eq:inversa-hensel-finita};
\item that a rational number should appear whose canonical expansion is not eventually
periodic, or an irrational number whose expansion is periodic;
\item that a finite deterministic emitter produces a tape that is
not eventually periodic;
\item that the catalog censuses do not sum to \(468\),
\(243\), and \(104\,976\), or that the measures of the three words do
not agree with the enumeration.
\end{enumerate}

The raw Hensel chart
\[
\mathbb Z_3^6\longrightarrow[0,1]
\]
is continuous and surjective; upon adjoining the integer part, its raw atlas covers \(\mathbb R\). This exact result does not imply that every raw fiber contains a representative that survives the additional filters. The published finite windows and the named readers of \(\pi,e,\varphi\) are verified within HMT, relative to their seeds, regions, orientations, and enriched rules.

\ifdefined\HMTInternalTraceability
For a filtered prefix, let \(T_u\) be the tree of compatible prolongations.
If \(T_u\) is finitely branching and has a node at every depth,
König's lemma yields an infinite branch; if, moreover, each stable level has
cardinality one, the branch is unique. This statement is conditional: the
universal nonemptiness of \(T_u\) for every prefix has not been proved. Thus, the
coverage of all fibers by filtered survivors and the well-defined descent of
native addition and multiplication to the Archimedean quotient remain an
\textsc{open programme}.
\fi
